% ===========================================================================
% Part X - Generative / Foundation Models: Large Language Models (Experiment 18)
% ===========================================================================

\part{Generative / Foundation Models: Large Language Models}

% ---------------------------------------------------------------------------
\chapter{Experiment 18: Large Language Model (LLM) Experiment}\label{exp:18}

\quickfacts
  {Dataset}{SMS Spam Collection (\texttt{data/sms\_spam.csv}, 5,572 messages; balanced subset 400 train / 200 test) + DistilBERT / DistilGPT2}
  {Key parameters}{Fine-tune: fresh 2-class head, 2 epochs, AdamW learning rate $5\times10^{-5}$, batch 16, max length 64}
  {Headline result}{Held-out accuracy 0.9600; spam F1 0.9592; greedy vs.\ nucleus generation demonstrated}

\section{Objective}
Demonstrate a reproducible LLM workflow for text classification and generation while
understanding tokenization, inference and evaluation.

\section{Suggested Dataset}
DistilBERT (SST-2) for tokenization and sentiment, DistilGPT2 for generation, and the real
SMS Spam Collection (\texttt{data/sms\_spam.csv}) for fine-tuning. A balanced subset
(300 ham + 300 spam) is split into 400 training and 200 stratified test messages.

\section{Required Libraries}
transformers, torch, pandas.

\section{Brief Theory}
Language models operate on \textbf{subword tokens} (Byte-Pair Encoding / WordPiece) mapped to
integer vocabulary IDs; an \textbf{attention mask} marks real tokens versus padding.
Self-attention computes
\[
\mathrm{Attention}(Q,K,V) = \mathrm{softmax}\!\left(\frac{QK^T}{\sqrt{d_k}} + \mathrm{mask}\right) V ,
\]
so every token attends to every other token. Causal generation predicts one token at a time
from $P(t_n \mid t_1, \dots, t_{n-1}) = \mathrm{softmax}(z_n / T)$: \textbf{greedy} decoding
takes the argmax each step, \textbf{temperature} $T$ scales the logits, and \textbf{top-$p$}
(nucleus) sampling restricts sampling to the smallest token set whose cumulative probability
exceeds $p$. \textbf{Fine-tuning} reuses pretrained weights, replaces the classification head
and trains briefly with a small learning rate.

\section{Procedure}
\begin{enumerate}
  \item Tokenize a sentence and inspect the tokens, IDs and attention mask.
  \item Run the sentiment pipeline on fixed test sentences.
  \item Generate text with greedy decoding and with temperature/top-$p$ sampling.
  \item Load the SMS Spam Collection and build a balanced subset (300 ham + 300 spam) with a
        400/200 stratified split.
  \item Fine-tune DistilBERT (fresh 2-class head, AdamW learning rate $5\times10^{-5}$,
        2 epochs, batch 16).
  \item Evaluate accuracy and F1 on the held-out messages and inspect example predictions.
\end{enumerate}

\section{Reference Implementation}
\begin{lstlisting}[style=mlcode,caption={Inference pipeline and the fine-tuning loop}]
classifier = pipeline("sentiment-analysis", model=model_id,
                      device=0 if torch.cuda.is_available() else -1)

ft_model = AutoModelForSequenceClassification.from_pretrained(model_id, num_labels=2,
                                                              ignore_mismatched_sizes=True)
optimizer = torch.optim.AdamW(ft_model.parameters(), lr=5e-5)
outputs = ft_model(input_ids=ids, attention_mask=mask, labels=labels)
outputs.loss.backward(); optimizer.step()
\end{lstlisting}

\section{Evaluation / Expected Output}
\begin{table}[htbp]
  \centering
  \caption{DistilBERT fine-tuned on the SMS Spam Collection (200 held-out messages).}
  \begin{tabular}{@{}lr@{}}
    \toprule
    \textbf{Item} & \textbf{Value} \\
    \midrule
    Fine-tune loss (epoch 1 $\rightarrow$ 2) & 0.705 $\rightarrow$ 0.079 \\
    Held-out accuracy & \best{0.9600} \\
    Held-out F1 (spam) & \best{0.9592} \\
    \bottomrule
  \end{tabular}
\end{table}

Example predictions (all correct): a prize-notification spam at 99.4\% confidence, a short ham
message at 99.5\%, and another spam at 89.7\%. Generation: greedy decoding produced the safe
deterministic continuation ``\textit{the way we think about the world}'', while nucleus
sampling ($T = 0.7$, $p = 0.9$) produced a different but coherent sentence --- the effect of
temperature and top-$p$.

The experiment demonstrates the complete practical pipeline and shows that fine-tuning a small
pretrained model on only 400 real messages already gives 96\% spam detection. Limitations: the
balanced subset ignores the natural 13\% spam prior of the full collection, the models are
small and can hallucinate, and prompts/data must be fixed for reproducibility.

\section{Viva / Discussion Questions}
\begin{description}
  \item[What is tokenization?] Splitting text into subword units that map to integer
  vocabulary IDs, letting a fixed vocabulary represent any word; special tokens and attention
  masks handle sequence framing and padding.
  \item[Pretraining vs fine-tuning vs prompting?] Pretraining learns general language
  representation from massive corpora; fine-tuning adapts those weights to a task with labeled
  data; prompting steers a frozen model through instructions or examples.
  \item[Why is LLM output not automatically ground truth?] The model generates the most
  plausible continuation, not verified facts; it can hallucinate, and confidence scores are not
  calibrated correctness probabilities.
  \item[How can prompt wording affect results?] Meaning, tone, format and label wording all
  change the conditional distribution the model samples from, so prompts must be fixed and
  reported.
\end{description}

\section{Complete Code Listing}
% ===========================================================================
% exp18_llm.tex - complete code listings for Experiment 18
% Source: notebooks/10_large_language_model.ipynb (executed)
% Repository: https://github.com/Atik203/ML-Algorithms
% ===========================================================================

\begin{codebox}[title={Core numerical and plotting libraries}]
# ---- Core numerical and plotting libraries ----
import os
import numpy as np
import pandas as pd
import matplotlib.pyplot as plt
import seaborn as sns

# ---- PyTorch: used for the fine-tuning demo ----
import torch
import torch.nn as nn
from torch.utils.data import DataLoader, Dataset
# ---- Dataset splitting + metrics for the fine-tune evaluation ----
from sklearn.model_selection import train_test_split
from sklearn.metrics import accuracy_score, f1_score
# ---- Hugging Face Transformers: tokenizers, models, pipelines, generation configs ----
from transformers import (
    AutoTokenizer,
    AutoModelForSequenceClassification,
    AutoModelForCausalLM,
    GenerationConfig,
    pipeline
)

# Styling setup (consistent look across all notebooks)
plt.style.use('seaborn-v0_8-whitegrid' if 'seaborn-v0_8-whitegrid' in plt.style.available else 'default')
plt.rcParams['font.sans-serif'] = 'DejaVu Sans'
plt.rcParams['figure.figsize'] = (10, 5)
plt.rcParams['figure.dpi'] = 110

torch.manual_seed(42)
torch.cuda.manual_seed_all(42)  # reproducible GPU training
device = torch.device('cuda' if torch.cuda.is_available() else 'cpu')  # GPU if available
print(f"PyTorch on {device}")
\end{codebox}

\begin{codebox}[title={Load a pre-trained DistilBERT tokenizer (fast, subword/WordPiece based)}]
# ---- Load a pre-trained DistilBERT tokenizer (fast, subword/WordPiece based) ----
model_id = "distilbert/distilbert-base-uncased-finetuned-sst-2-english"
print(f"Loading Tokenizer: {model_id}...")
tokenizer = AutoTokenizer.from_pretrained(model_id)

sample_text = "Machine learning and deep neural networks are transforming modern science!"
# Tokenize with padding/truncation so the output has a fixed shape (here 20 tokens)
encoded = tokenizer(
    sample_text,
    padding="max_length",
    max_length=20,
    truncation=True,
    return_tensors="pt"
)

# Convert integer IDs back to readable subword tokens
tokens = tokenizer.convert_ids_to_tokens(encoded['input_ids'][0])
print(f"Original Text: {sample_text}\n")
print(f"Subword Tokens ({len(tokens)}):\n", tokens)
print(f"Input IDs:      \n", encoded['input_ids'][0].tolist())
print(f"Attention Mask: \n", encoded['attention_mask'][0].tolist())
\end{codebox}

\begin{codebox}[title={Load a pre-trained sentiment-analysis pipeline (DistilBERT SST-2)}]
# ---- Load a pre-trained sentiment-analysis pipeline (DistilBERT SST-2) ----
classifier = pipeline("sentiment-analysis", model=model_id, device=0 if torch.cuda.is_available() else -1)

test_sentences = [
    "The new neural network model demonstrated extraordinary accuracy and remarkable speed.",
    "The dataset was poorly labeled, noisy, and the model completely failed to converge.",
    "The gradient descent algorithm performed reasonably well, meeting standard expectations.",
    "A catastrophic failure occurred during training due to exploding gradient issues."
]

# One forward pass per sentence -> label + calibrated confidence score
predictions = classifier(test_sentences)

# Collect results in a table for display
results_table = []
for sent, pred in zip(test_sentences, predictions):
    results_table.append({
        "Sentence": sent,
        "Predicted Label": pred['label'],
        "Confidence Score": f"{pred['score']*100:.2f}%"
    })

df_preds = pd.DataFrame(results_table)
display(df_preds)
\end{codebox}

\begin{codebox}[title={Load a small causal LM (DistilGPT2) for text generation}]
# ---- Load a small causal LM (DistilGPT2) for text generation ----
gen_model_id = "distilbert/distilgpt2"
print(f"Loading Text Generator: {gen_model_id}...")
# clean_up_tokenization_spaces=False keeps BPE spacing intact
generator = pipeline("text-generation", model=gen_model_id,
                     device=0 if torch.cuda.is_available() else -1, clean_up_tokenization_spaces=False)

prompt = "Artificial Intelligence will fundamentally transform"

# 1. Greedy Search: always pick the most likely next token (deterministic)
greedy_output = generator(
    prompt, generation_config=GenerationConfig(max_new_tokens=40, do_sample=False)
)[0]['generated_text']

# 2. Temperature + Top-p (Nucleus) Sampling: sample from the most likely tokens up to cumulative p
sampled_output = generator(
    prompt, generation_config=GenerationConfig(max_new_tokens=40, do_sample=True, temperature=0.7, top_p=0.9)
)[0]['generated_text']

print(f"PROMPT: '{prompt}'\n")
print(f"--- 1. GREEDY SEARCH OUTPUT ---\n{greedy_output.strip()}\n")
print(f"--- 2. NUCLEUS SAMPLING (T=0.7, p=0.9) OUTPUT ---\n{sampled_output.strip()}")
\end{codebox}

\begin{codebox}[title={Real dataset: SMS Spam Collection (5,572 labeled messages)}]
# ---- Real dataset: SMS Spam Collection (5,572 labeled messages) ----
sms = pd.read_csv('../data/sms_spam.csv')
sms['y'] = (sms['label'] == 'spam').astype(int)
print(f"SMS Spam Collection: {len(sms)} messages | ham {(sms.y == 0).sum()}, spam {(sms.y == 1).sum()}")

# Balanced subset for a fast fine-tune demo: 300 ham + 300 spam
ham = sms[sms.y == 0].sample(n=300, random_state=42)
spam = sms[sms.y == 1].sample(n=300, random_state=42)
subset = pd.concat([ham, spam]).sample(frac=1, random_state=42).reset_index(drop=True)

# Stratified split keeps the ham/spam ratio identical in train and test
train_df, test_df = train_test_split(subset, test_size=200, random_state=42, stratify=subset.y)
print(f"Train messages: {len(train_df)} | Test messages: {len(test_df)} | spam share: {subset.y.mean():.1%}")

class SMSDataset(Dataset):
    """Tokenize SMS texts and expose them as PyTorch samples."""
    def __init__(self, texts, labels, tokenizer):
        self.encodings = tokenizer(list(texts), padding=True, truncation=True, max_length=64, return_tensors="pt")
        self.labels = list(labels)

    def __getitem__(self, idx):
        item = {key: val[idx] for key, val in self.encodings.items()}
        item['labels'] = torch.tensor(self.labels[idx], dtype=torch.long)
        return item

    def __len__(self):
        return len(self.labels)

train_loader = DataLoader(SMSDataset(train_df.text, train_df.y, tokenizer), batch_size=16, shuffle=True)
test_loader = DataLoader(SMSDataset(test_df.text, test_df.y, tokenizer), batch_size=32)

# ---- Load the pre-trained classifier with a fresh 2-class head ----
# ignore_mismatched_sizes allows replacing the original SST-2 head with a new one.
ft_model = AutoModelForSequenceClassification.from_pretrained(model_id, num_labels=2, ignore_mismatched_sizes=True)
ft_model.to(device)
optimizer = torch.optim.AdamW(ft_model.parameters(), lr=5e-5)  # small LR for fine-tuning

print("Pre-trained Transformer loaded for fine-tuning!")
\end{codebox}

\begin{codebox}[title={Fine-tuning loop (2 epochs on 400 real SMS messages)}]
# ---- Fine-tuning loop (2 epochs on 400 real SMS messages) ----
epochs = 2
ft_model.train()
epoch_losses = []

for epoch in range(epochs):
    running_loss = 0.0
    for batch in train_loader:
        optimizer.zero_grad()

        # Move the batch to the selected device
        input_ids = batch['input_ids'].to(device)
        attention_mask = batch['attention_mask'].to(device)
        labels = batch['labels'].to(device)

        # Forward pass: the model computes the classification loss internally
        outputs = ft_model(input_ids=input_ids, attention_mask=attention_mask, labels=labels)
        loss = outputs.loss
        loss.backward()      # backpropagate through the whole transformer
        optimizer.step()     # update weights

        running_loss += loss.item() * len(input_ids)   # de-average the batch loss

    avg_loss = running_loss / len(train_df)
    epoch_losses.append(avg_loss)
    print(f"Epoch [{epoch+1}/{epochs}] Fine-Tuning Loss: {avg_loss:.4f}")
\end{codebox}

\begin{codebox}[title={Evaluate the fine-tuned model on held-out real SMS messages}]
# ---- Evaluate the fine-tuned model on held-out real SMS messages ----
ft_model.eval()
preds, targets = [], []

with torch.no_grad():
    for batch in test_loader:
        # Move the batch to the same device as the fine-tuned model (GPU when available)
        input_ids = batch['input_ids'].to(device)
        attention_mask = batch['attention_mask'].to(device)
        logits = ft_model(input_ids=input_ids, attention_mask=attention_mask).logits
        preds += logits.argmax(-1).tolist()       # predicted class (0 = ham, 1 = spam)
        targets += batch['labels'].tolist()

print(f"Held-out test accuracy: {accuracy_score(targets, preds):.4f}")
print(f"Held-out test F1 (spam): {f1_score(targets, preds):.4f}")

# ---- Show a few individual predictions with confidence ----
labels_map = {0: "ham", 1: "spam"}
with torch.no_grad():
    for i in range(3):
        message = test_df.text.iloc[i]
        enc = tokenizer(message, truncation=True, max_length=64, return_tensors="pt").to(device)
        probs = torch.softmax(ft_model(**enc).logits, dim=1).cpu().numpy()[0]
        print(f"\nMessage: {message[:75]}...")
        print(f"True: {labels_map[test_df.y.iloc[i]]} | Predicted: {labels_map[probs.argmax()]} ({probs.max()*100:.1f}%)")
\end{codebox}

